\documentclass[conference]{IEEEtran}
\usepackage[utf8]{inputenc}
\usepackage[T1]{fontenc}
\usepackage{cite}
\usepackage{amsmath,amssymb}
\usepackage{graphicx}
\usepackage{booktabs}
\usepackage{url}

\title{Verifier‑Grounded Evaluation of LLM‑Generated Network Configurations: A Preregistered NetConfig‑Semantic Study}

\author{
\IEEEauthorblockN{Autonomous AI Researcher}
\IEEEauthorblockA{CNSM 2026 Agentic AI Researcher Track}
}

\begin{document}

\maketitle

\begin{abstract}
We report a fully preregistered, artifact‑anchored paired experiment (CAND‑2026‑001‑NetConfig‑Semantic‑v1; preregistration SHA‑256 7db1663cf3982c8ea1e16c3dd7c0348f356bb4cbe3978ecc75bb59fdedb71a37) comparing a deterministic template baseline to a single‑pass LLM generator (declared\_model\_version = gpt‑5‑mini) on N = 120 synthetic intent\ensuremath{\rightarrow}configuration tasks using a deterministic validator. Under the frozen confirmatory semantics (shared\_initial\_candidate per pair) all confirmatory episodes completed: baseline = 120/120 verifier passes; guarded (gpt‑5‑mini, seed\_0) = 120/120 verifier passes. Paired difference = 0.00; 95\% bootstrap percentile CI [0.00, 0.00] (10,000 resamples, seed = 1729); exact McNemar two‑sided p = 1.0. We (i) publish the exact execution and analysis artifacts and SHA‑256 fingerprints needed to reproduce the run (see execution/execution\_manifest.json in the artifact bundle), (ii) diagnose---using archived artifacts---why the paired outcome is uninformative about independent generative reliability (preregistered shared\_initial\_candidate design, template‑tight task synthesis, and validator--task coupling), and (iii) provide artifact‑grounded, runnable modifications (independent\_per\_condition sampling, multi‑seed confirmatory design, validator diversity, and expanded task heterogeneity) that preserve reproducibility while producing informative independent comparisons. The immutable initial master prompt is archived (provenance/master\_prompt.txt; SHA‑256 1872df1e1805d2d96940456ca016bd665d1d5196add77f5acdf1582bb39b15ba).
\end{abstract}

\section{Introduction}
Motivation. Translating operator intent into device configuration is a core NetOps task where correctness is safety‑critical: incorrect commands can break reachability, violate ACLs, or create unsafe transient states during multi‑step changes. Deterministic offline verifiers (reachability/ACL/routing analyzers such as Batfish‑class tools) are widely used to validate intended changes before deployment and provide an operational oracle for generated configuration artifacts [https://doi.org/10.1145/3603269.3604866]. Large language models (LLMs) offer rapid intent\ensuremath{\rightarrow}config generation but raise reliability questions: how often do independently sampled LLM candidates satisfy operational verifier constraints? How do LLM pipelines compare with deterministic template baselines when correctness is judged by a deterministic validator?

Research question and contribution. After a fresh autonomous literature and gap analysis we preregistered (CAND‑2026‑001‑NetConfig‑Semantic‑v1) a paired confirmatory test asking: does a single‑pass LLM generator (gpt‑5‑mini) have a lower verifier pass rate than a deterministic template baseline across N = 120 intent\ensuremath{\rightarrow}config tasks? Our primary contribution is an artifact‑anchored, fully reproducible preregistered execution + analysis bundle and a transparent diagnosis of why the preregistered confirmatory semantics (shared\_initial\_candidate) produced a degenerate but correct paired outcome. We (1) publish the exact artifacts and SHA‑256 fingerprints required for byte‑for‑byte reruns; (2) report confirmatory counts and preregistered statistical inferences grounded in the archived traces (baseline = 120/120, guarded = 120/120; paired difference = 0.00; bootstrap CI [0.00,0.00]; McNemar p = 1.0); and (3) provide artifact‑grounded, immediately runnable experimental modifications that preserve reproducibility while answering the operational question of independent generative reliability.

\section{Related Work}
Verifier‑grounded NetOps and intent\ensuremath{\rightarrow}configuration. Offline semantics analyzers that compute network final and transient state from device configurations are established engineering controls in pre‑deployment validation. Batfish and related tools demonstrate how high‑fidelity configuration analysis enables operators to detect semantic errors before changes reach production; Batfish's engineering lessons emphasize scalability, fidelity, and the operational value of deterministic verification in NetOps workflows [10.1145/3603269.3604866]. These verifiers produce machine‑checkable oracles (reachability matrices, ACL equivalence, route‑preservation checks and transient‑state invariants) that directly map to concrete operational failure modes (lost reachability, unintended route leaks, unsafe transient downtimes).

LLM\ensuremath{\rightarrow}NetOps evaluations and generator--validator patterns. Recent work on intent‑driven configuration generation and related pipelines has converged on generate--validate--repair architectural patterns: an LLM proposes candidate configurations, a deterministic verifier checks semantics, and a repair loop attempts to fix validator failures. Surveys and position papers on LLM‑based network management identify this pattern and emphasize that deterministic validators are necessary to move beyond syntactic or fluency‑based evaluation metrics [10.1002/nem.70029]. Empirical studies of LLM repair loops in adjacent domains show diminishing returns from unbounded iteration and highlight that evaluation outcomes depend strongly on orchestration and feedback design rather than model identity alone; these findings motivate careful treatment of repair budgets and sampling semantics when designing experiments that measure operational reliability.

Evaluation pitfalls: semantic scoring vs. generative variability. Prior evaluations have sometimes conflated different scientific questions: (i) how often a particular single candidate produced by a pipeline passes a verifier (a candidate‑level success), (ii) how reliably independently sampled candidates satisfy verifier constraints (a generative‑reliability probability), and (iii) whether a repair loop can raise a failed candidate to pass. Studies that report per‑candidate pass rates without specifying generation semantics or seed treatment risk ambiguity. Our preregistered design explicitly fixed generation semantics (shared\_initial\_candidate) and thus answered a narrowly specified estimand: whether a deterministic verifier treats identical artifacts differently across scoring conditions---an important but distinct question from independent LLM reliability. The distinction maps directly onto operational decisions: operators choosing to accept a single canonical LLM candidate with verification differ from operators who rely on repeated sampling or repair loops to achieve probabilistic guarantees.

Validator--task coupling and construct validity. When tasks are programmatically templated and validators enforce the very workflow\_policies encoded in those templates, canonical candidates that mirror the template structure will tend to pass deterministically. This validator--task coupling can inflate pass rates relative to open distributions of natural intents. The literature on intent‑driven generators recognizes this trade‑off: narrow, template‑aligned tasks allow reproducible, deterministic evaluation useful for principled method comparisons, but they limit external validity to more heterogeneous operational intents (surveyed in [10.1002/nem.70029]).

Preregistration, artifact release, and reproducibility. A complementary body of work emphasizes transparent artifact publication (model traces, seeds, verifier logs) and preregistration to make explicit what a given test measures and to permit deterministic reruns under alternative semantics. Our contribution adopts this practice: we preregistered the estimand, generation semantics, and analysis plan, and we release the full per‑episode scoring traces and SHA‑256 manifest so that others can reexecute the experiment under different sampling semantics (for example, independent\_per\_condition sampling or multi‑seed confirmatory designs) while preserving deterministic reproducibility of each variant. In net effect, verifier grounding plus exhaustive artifact publication clarifies both the scientific question answered and the minimal modifications required to answer adjacent operational questions.

\section{Methodology}
Preregistration and frozen design. We preregistered study CAND‑2026‑001‑NetConfig‑Semantic‑v1 prior to observing any model outputs; the preregistration artifact is archived (SHA‑256 7db1663cf3982c8ea1e16c3dd7c0348f356bb4cbe3978ecc75bb59fdedb71a37). The preregistration specified: paired confirmatory design across N = 120 tasks (40 tasks per topology class: edge, datacenter, multi‑AS), primary estimand paired\_success\_rate\_difference\_guarded\_minus\_baseline, confirmatory generation\_semantics = "shared\_initial\_candidate", initial canonical seed\_0 per pair, validator = deterministic\_netops\_validator\_v5, and the exact analysis plan (paired\_binary\_analysis\_v1) including bootstrap parameters (10,000 resamples, seed = 1729). The full execution manifest and per‑file SHA‑256s are published (execution/execution\_manifest.json; see artifact manifest in submission bundle).

Task synthesis and templates. Tasks were programmatically synthesized from a deterministic template bank (generator\_id deterministic\_netops\_task\_generator\_v5, version 5.0). Each task manifest entry encodes: initial\_state, expected\_final\_state, an explicit required\_sequence (the minimal change sequence that must be executed to satisfy intent safely), per‑task workflow\_policies (e.g., must\_be\_down\_for\_fields, minimum\_active\_in\_groups), allowed\_touched\_fields, and a difficulty annotation (changed\_interface\_count, transient\_rule\_count). Representative task manifest entries and their SHA‑256s are archived (e.g., execution/task\_manifest.jsonl; see task examples under artifact\_examples). Templates intentionally encode operationally meaningful transient constraints (e.g., make‑before‑break switchover, atomic MTU changes) so that verifier checks carry operational semantics rather than purely syntactic checks.

Model, orchestration, and generation semantics. The declared LLM was gpt‑5‑mini accessed via the hosted adapter family hosted\_netops\_gvr\_v1 (execution/model\_configuration.json; SHA‑256 a40e96e212ab87da251ac0d6dbb75bc543afa13438b5a6b9ebd073597ffdd820). The preregistered confirmatory generation semantics produced a single canonical candidate per pair (shared\_initial\_candidate) and reused that identical candidate for both baseline and guarded scoring episodes; the canonical candidate files are archived at execution/responses/*-shared-initial.txt (representative SHA‑256s: task‑000001 shared initial SHA‑256 8f38c8becd7e1e36..., task‑000002 SHA‑256 ae967f70dfb6ccb8..., task‑000003 SHA‑256 f7f4db249ecbbce...).

Deterministic validator and scoring. Each candidate was scored by deterministic\_netops\_validator\_v5 (validator\_version 5.0). The validator parses commands, normalizes configuration sequences, enforces per‑task workflow\_policies, runs transient safety checks (e.g., minimum\_active\_in\_groups during the change), and computes final\_state reachability and ACL invariants. The scoring rule used for the confirmatory test was full\_intent\_constraint\_validation: a binary pass (VALID\_CONFIGURATION = 1) requires that all required\_commands and transient constraints are satisfied. Each scoring episode emits a detailed JSON trace archived at execution/scoring/*.json (representative: execution/scoring/task-000001-guarded.json SHA‑256 0d56c1add7c5a57f...).

Execution accounting, retries, contamination heuristics. Planned\_episode\_count = 240 (120 pairs \ensuremath{\times} 2 conditions). Completed\_episode\_count = 240; model\_calls\_used = 120 (shared\_initial generation per pair counted once); failed\_episode\_count = 0. A preregistered retry policy allowed up to two automated retries for infrastructure failures; none were required. Contamination heuristics (exact‑match and normalized n‑gram overlap) were run as preregistered and recorded (analysis/contamination\_summary.csv). All provider call traces and call\_cache artifacts are included (execution/provider\_calls/*.json, execution/call\_cache/*) to permit external exposure analysis.

\section{Results}
Confirmatory counts and inferential summary (artifact grounding). All confirmatory scoring episodes completed and are traceable in the execution manifest (execution/execution\_manifest.json). The archived confirmatory counts are: baseline\_success\_count = 120/120; guarded\_success\_count = 120/120; complete\_pair\_count = 120. The paired contingency table is degenerate: n\_11 = 120, n\_10 = 0, n\_01 = 0, n\_00 = 0 (analysis/paired\_contingency\_table.csv SHA‑256 9f25790c7eeaafb3...). The primary estimator (mean of guarded\_i - baseline\_i) = 0.00. The preregistered bootstrap (10,000 resamples, seed 1729) yields a 95\% percentile CI [0.00, 0.00]; exact McNemar two‑sided p = 1.0. Analysis artifacts (condition\_summary.csv SHA‑256 9c77c96f37c4b6ff..., paired\_difference\_figure.svg SHA‑256 ae5b8e0c644f8cf7...) are published in the bundle.

Representative validator traces (artifact examples). The per‑episode JSON traces record parsed\_command\_count, required\_command\_count, normalized\_configuration, validation\_before and validation\_after final\_state dictionaries, and violation\_count. Examples: execution/scoring/task-000001-guarded.json (SHA‑256 0d56c1add7c5a57f...) shows parsed\_command\_count = 4, required\_command\_count = 4, violation\_count = 0 and contains explicit before/after final\_state mappings. The canonical candidate used for that pair is archived at execution/responses/task-000001-shared-initial.txt (SHA‑256 8f38c8becd7e1e36...). Multiple pairs exhibit identical baseline/guarded/shared response fingerprints (e.g., task‑000003 artifacts all share SHA‑256 f7f4db249ecbbce...), demonstrating that identical candidates were evaluated across both conditions.

Why the confirmatory outcome is degenerate (artifact‑grounded diagnosis). The degenerate contingency (all pairs n\_11) follows directly from two preregistered, artifact‑visible design choices: (1) the confirmatory generation\_semantics set to shared\_initial\_candidate caused a single canonical candidate to be reused for both baseline and guarded episodes per pair (see execution/responses/*-shared-initial.txt SHA‑256s), and (2) the task templates encoded explicit required\_sequences and workflow\_policies that canonical candidates satisfied. Under these conditions a deterministic validator necessarily returns identical pass/fail judgements for identical inputs, hence the paired difference must be zero. The execution and scoring artifacts (execution/responses/* and execution/scoring/*) fully document this chain of causation and are published to enable exact reruns.

Contamination and exposure. Preregistered contamination heuristics flagged no confirmatory items under the chosen thresholds (analysis/contamination\_summary.csv). All provider call traces and call\_cache entries are published to permit stronger external exposure checks; absence of a flag in the preregistered heuristic is not proof of zero pretraining exposure and is explicitly noted in the Disclosure.

\section{Discussion}
Interpretation of confirmatory inference. The reported confirmatory statistics are correct for the preregistered estimand: they test whether a deterministic verifier treats identical artifacts differently under two scoring conditions. They do not, however, estimate the operational quantity of independent LLM generative reliability (the probability that independently sampled LLM candidates satisfy verifier constraints). For deployment decisions operators need independent\_per\_condition pass probabilities; the archived artifacts permit reruns under those semantics without creating new task definitions.

Artifact‑grounded remediation and runnable next experiments. Using the published bundle one can immediately perform informative reruns while preserving reproducibility: (A) independent\_per\_condition confirmatory sampling --- generate separate seeds per condition per task (e.g., seed\_0 baseline, seed\_1 guarded), rescore with deterministic\_netops\_validator\_v5 and compute paired contrasts; (B) multi‑seed confirmatory design --- draw K >= 3 independent seeds per condition per task and estimate per‑condition pass probabilities with bootstrap CIs; (C) validator diversification --- add a second independent verifier (e.g., reachability emulator) and report intersection/union pass rates to reduce validator--task coupling. Because the generator, provider traces, and scoring harness are fully archived (execution/model\_configuration.json SHA‑256 a40e96e2..., execution/provider\_calls/*.json, execution/scoring/*.json), these reruns are deterministic and reproducible.

Threats to validity (artifact‑identified). Internal validity: the shared\_initial design intentionally removed generation variance and therefore precluded testing independent generative reliability. Construct validity: template‑aligned tasks and a single deterministic validator increase the chance that canonical candidates satisfy required\_sequences, possibly overstating success in open distributions. External validity: a single declared model (gpt‑5‑mini) and synthetic tasks limit generalization to other model families, open natural language intents, and production topologies. Conclusion validity: a degenerate contingency table produces trivial but correct inferential outputs; interpreting them beyond their precise estimand would be inappropriate. All these threats are documented in the published traces.

Operational implications for NetOps practice. Our artifact‑anchored diagnosis shows that (1) verifier grounding is necessary but not sufficient---experimental semantics must match the deployment question; (2) transparent preregistration plus exhaustive artifact publication materially clarifies what a given confirmatory test measures; and (3) simple rerun protocols (independent sampling, multi‑seed replication, validator diversity) convert a degenerate confirmatory design into an informative operational comparison without sacrificing reproducibility.

\section{Limitations}
\begin{itemize}
\item Controlled synthetic tasks: programmatic templates produced narrow required\_sequences and workflow\_policies that favour canonical solutions and limit external generalizability to heterogeneous operational intents.
\item Confirmatory shared\_initial semantics: the preregistered generation\_semantics = 'shared\_initial\_candidate' supplied identical canonical candidates to both conditions per pair, reducing independence between conditions and causing the degenerate paired outcome.
\item Single canonical seed for confirmatory test: the primary analysis used one canonical seed per task; preregistered exploratory multi‑seed analyses (seeds\_1..4) were not included in the confirmatory inference reported here.
\item Validator--task coupling: the deterministic validator enforces constraints encoded in the task templates, which can increase pass rates for template‑aligned candidates and reduce construct validity for open distributions.
\item Possible training‑data exposure: preregistered contamination heuristics found no flags under chosen thresholds, but absence of a flag is not proof of zero exposure to related artifacts in model pretraining.
\item Single‑model, single‑validator run: results apply only to gpt‑5‑mini under the recorded configuration and deterministic\_netops\_validator\_v5; they do not generalize across model families, agentic pipelines, or other validators.
\end{itemize}

\section{Conclusion}
We executed a fully preregistered, verifier‑grounded paired experiment and released a complete artifact bundle enabling byte‑level reproduction (execution/execution\_manifest.json and analysis/*). Under the preregistered confirmatory semantics (shared\_initial\_candidate) both the deterministic template baseline and the gpt‑5‑mini guarded condition validated on all N = 120 tasks (both 120/120). Archived evidence (shared\_initial candidate files and validator scoring traces) demonstrates that identical canonical candidates per pair and template‑tight task constraints caused the degenerate paired outcome; consequently the confirmatory paired result does not provide evidence about independent LLM generative reliability in operational settings. We provide artifact‑grounded, runnable modifications (independent\_per\_condition sampling, multi‑seed confirmatory execution, validator diversification, task heterogeneity expansion) that preserve reproducibility and enable informative independent comparisons. All execution and analysis artifacts---including the immutable master prompt reference (provenance/master\_prompt.txt; SHA‑256 1872df1e1805d2d96940456ca016bd665d1d5196add77f5acdf1582bb39b15ba) and the preregistration SHA‑256---are included in the submission bundle to permit exact audit and follow‑up experiments. Transparent preregistration combined with exhaustive artifact publication clarifies precisely what a confirmatory test measures and accelerates corrective reruns that answer the operational questions NetOps practitioners require for deployment decisions.

\section*{Disclosure Statement}
Mandatory disclosure (CNSM Agentic AI Researcher track). LLM(s) and provider: declared\_model\_version = gpt-5-mini accessed via provider adapter 'openai\_responses' and adapter family 'hosted\_netops\_gvr\_v1' (execution/model\_configuration.json SHA‑256 a40e96e212ab87da251ac0d6dbb75bc543afa13438b5a6b9ebd073597ffdd820). Orchestration/framework: hosted\_netops\_gvr\_v1 executed the preregistered pipeline; deterministic scoring used deterministic\_netops\_validator\_v5 (validator\_version 5.0). Preregistration: CAND-2026-001-NetConfig-Semantic-v1 (frozen prior to execution); preregistration SHA‑256 7db1663cf3982c8ea1e16c3dd7c0348f356bb4cbe3978ecc75bb59fdedb71a37. Exact initial master prompt: preserved as an immutable archived file provenance/master\_prompt.txt (SHA‑256 1872df1e1805d2d96940456ca016bd665d1d5196add77f5acdf1582bb39b15ba). Human role and autonomy boundary: the Research Director selected the topic family (Generative AI and LLMs for NetOps) and approved the initial master prompt before the autonomy lock. After the autonomy lock the autonomous pipeline performed literature review, research‑gap selection, preregistration, code generation, experiment execution, deterministic validation, statistical analysis, autonomous internal peer review, manuscript drafting and revision. No human manually edited technical manuscript content after the autonomy lock. Preregistered confirmatory generation semantics and consequence: confirmatory generation\_semantics was set to 'shared\_initial\_candidate' so each pair received an identical canonical candidate for both baseline and guarded episodes; representative shared candidate artifacts: execution/responses/task-000001-shared-initial.txt (SHA‑256 8f38c8becd7e1e36...), execution/responses/task-000002-shared-initial.txt (SHA‑256 ae967f70dfb6c...), execution/responses/task-000003-shared-initial.txt (SHA‑256 f7f4db249ecbbce...). Validator and scoring: deterministic\_netops\_validator\_v5 performed normalized parsing, intent constraint checking, transient safety checks, and emitted per‑episode JSON scoring traces archived under execution/scoring/*.json (representative SHA‑256s: execution/scoring/task-000001-guarded.json SHA‑256 0d56c1add7c5a57f...). Execution accounting and retries: planned\_episode\_count = 240, completed\_episode\_count = 240, model\_calls\_used = 120, failed\_episode\_count = 0; preregistered retry policy allowed up to 2 automatic retries for infrastructure failures; none were required. Contamination checks and reproducibility: preregistered string‑overlap heuristics found no confirmatory flags under chosen thresholds, but absence of a flag is not proof of zero exposure to related artifacts in model pretraining. All artifacts, seeds, code, and per‑file SHA‑256 hashes required for exact reproduction are released in the submission bundle (execution/execution\_manifest.json contains the exhaustive manifest). The initial master prompt is referenced immutably by path and SHA‑256 above.

\begin{thebibliography}{99}
\bibitem{ref1} https://doi.org/10.1145/3603269.3604866
\bibitem{ref2} https://doi.org/10.23919/cnsm62983.2024.10814407
\bibitem{ref3} https://doi.org/10.1002/nem.70029
\end{thebibliography}

\end{document}
